% Book 3, Year 1, chapter 26 -- minimum spanning trees and union--find.
% Provenance (never printed): sources/book3/BRIEFS.md, chapter 26.
% ledger: b1:kruskal-1956, b1:prim-1957, b1:tarjan-uf-1975
\chapter{Minimum Spanning Trees and Union--Find}\label{ch:b1:spanning-trees}

Brenmoor wants fibre to every place, laid along its roads. Every place must be connected to every other, through the
cable if not directly, and cable costs by the metre. The town has 61 kilometres of road; how little of it can carry the
fibre? The cheapest network never closes a loop, since removing one road of a loop disconnects nothing: it is a tree,
and the question is to find the lightest one. Two \omterm{def:b1:greedy-proofs:greedy}{greedy algorithms} find it, and their proofs rest on one \omterm{met:b1:greedy-proofs:exchange}{exchange
argument}. The faster of them needs to answer, again and again, ``are these two places already connected?'', and the
data structure that answers it in almost constant time is the second subject of the chapter.

\begin{recall}
A \omterm{def:b1:graph-traversals:free-tree}{free tree} on $n$ \omterm{def:b1:graph-traversals:graph}{vertices} is connected with $n - 1$ edges (\cref{prop:b1:graph-traversals:free-tree}). \omterm{def:b1:greedy-proofs:greedy}{Greedy
algorithms} are proved by exchange (\cref{ch:b1:greedy-proofs}); \omterm{def:b1:heaps:binary-heap}{binary heaps} are in \cref{ch:b1:heaps}.
\end{recall}

\section{Spanning trees and the cut property}\label{sec:b1:spanning-trees:cut}

\begin{definition}[Spanning tree, minimum spanning tree]\label{def:b1:spanning-trees:mst}
A \emph{spanning tree}\index{spanning tree} of a connected undirected \omterm{def:b1:graph-traversals:graph}{graph} $G$ is a subgraph that contains every \omterm{def:b1:graph-traversals:graph}{vertex}
of $G$ and is a \omterm{def:b1:graph-traversals:free-tree}{free tree}. In a \omterm{def:b1:graph-traversals:weighted}{weighted graph}, its weight is the sum of the weights of its edges, and a \emph{minimum
spanning tree}\index{minimum spanning tree} is a spanning tree of smallest weight.
\end{definition}

\begin{definition}[Cut, crossing edge]\label{def:b1:spanning-trees:cut}
A \emph{cut}\index{cut (of a graph)} of a \omterm{def:b1:graph-traversals:graph}{graph} is a partition of its \omterm{def:b1:graph-traversals:graph}{vertices} into two non-empty sets $S$ and $V
\setminus S$. An edge \emph{crosses}\index{crossing edge} the cut if it has one end in each set.
\end{definition}

\begin{theorem}[Cut property, cycle property]\label{thm:b1:spanning-trees:cut-property}
Let $G$ be a connected \omterm{def:b1:graph-traversals:weighted}{weighted graph}. (\omterm{def:b1:spanning-trees:cut}{Cut} property) For any \omterm{def:b1:spanning-trees:cut}{cut}, a lightest edge crossing it belongs to some \omterm{def:b1:spanning-trees:mst}{minimum
spanning tree}; if it is the only lightest one, to every \omterm{def:b1:spanning-trees:mst}{minimum spanning tree}. (Cycle property) An edge strictly heavier
than every other edge of some cycle belongs to no \omterm{def:b1:spanning-trees:mst}{minimum spanning tree}.
\end{theorem}

\begin{proof}
Let $e = uv$ be a lightest edge crossing the \omterm{def:b1:spanning-trees:cut}{cut} $(S, V \setminus S)$, $u \in S$, and $T$ a \omterm{def:b1:spanning-trees:mst}{minimum spanning tree} without
$e$. The path from $u$ to $v$ in $T$ starts in $S$ and ends outside, so one of its edges, $f$, \omterm{def:b1:spanning-trees:cut}{crosses} the \omterm{def:b1:spanning-trees:cut}{cut}. Removing
$f$ from $T$ splits it into two trees, one containing $u$ and one $v$; adding $e$ joins them again. The result $T' = T
- f + e$ is a \omterm{def:b1:spanning-trees:mst}{spanning tree} of weight $w(T) - w(f) + w(e) \leq w(T)$, so it is minimum and contains $e$. If $e$ is the
only lightest edge, $w(e) < w(f)$ and $T'$ would be lighter than $T$: so no \omterm{def:b1:spanning-trees:mst}{minimum spanning tree} lacks $e$. For the
cycle property, let $e$ be the heaviest edge of a cycle $C$, in a \omterm{def:b1:spanning-trees:mst}{spanning tree} $T$. Removing $e$ splits $T$ into two
parts, a \omterm{def:b1:spanning-trees:cut}{cut} that $e$ \omterm{def:b1:spanning-trees:cut}{crosses}; the rest of $C$ joins the ends of $e$, so it has another edge $f$ crossing that \omterm{def:b1:spanning-trees:cut}{cut}, and
$T - e + f$ is a lighter \omterm{def:b1:spanning-trees:mst}{spanning tree}.
\end{proof}

\begin{omfigure}
\begin{tikzpicture}[v/.style={circle, draw, fill=white, minimum size=4mm, inner sep=0pt}, font=\scriptsize]
\fill[omDef!12, rounded corners=4mm] (-0.6,-1.6) rectangle (2.4,1.6);
\node[font=\small, omDef] at (0.9,1.9) {$S$}; \node[font=\small] at (5.1,1.9) {$V \setminus S$};
\node[v] (a) at (0,0.9) {}; \node[v] (u) at (1.6,0.2) {}; \node[v] (b) at (0.3,-1.0) {};
\node[v] (v) at (4.0,0.5) {}; \node[v] (c) at (5.6,1.0) {}; \node[v] (d) at (5.2,-0.9) {}; \node[v] (x) at (3.9,-1.1)
{};
\draw[very thick] (a) -- (u) -- (b); \draw[very thick] (v) -- (c) -- (d) -- (x);
\draw[very thick, omThm] (b) -- node[below] {$f$, in $T$} (x);
\draw[thick, dashed, omDef] (u) -- node[above] {$e$} (v);
\node[anchor=south east] at (u.north west) {$u$}; \node[anchor=south west] at (v.north east) {$v$};
\node[anchor=west, align=left] at (6.2,0.2) {$e$: a lightest crossing edge\\ $T - f + e$ is a spanning tree\\ no heavier than $T$};
\end{tikzpicture}
\par\smallskip{\small The exchange in the \omterm{def:b1:spanning-trees:cut}{cut} property: the tree $T$ (thick) joins $u$ to $v$ by a path that \omterm{def:b1:spanning-trees:cut}{crosses} the
\omterm{def:b1:spanning-trees:cut}{cut} at some edge $f$; swapping $f$ for the lightest \omterm{def:b1:spanning-trees:cut}{crossing edge} $e$ keeps a \omterm{def:b1:spanning-trees:mst}{spanning tree} and cannot add weight.}
\end{omfigure}

With distinct weights the \omterm{def:b1:spanning-trees:mst}{minimum spanning tree} is unique (\cref{exo:b1:spanning-trees:10}): every edge of it is the
only lightest edge across some \omterm{def:b1:spanning-trees:cut}{cut}.

\section{Kruskal's algorithm}\label{sec:b1:spanning-trees:kruskal}

\begin{method}[Kruskal's algorithm]\label{met:b1:spanning-trees:kruskal}
Take the edges by increasing weight; keep an edge if it joins two different trees of the \omterm{def:b1:graph-traversals:free-tree}{forest} kept so far, that is,
if it closes no cycle. Stop at $n - 1$ edges.
\end{method}

Each kept edge $uv$ is a lightest edge crossing the \omterm{def:b1:spanning-trees:cut}{cut} between the tree of $u$ and the rest, since every lighter edge
crossing it was examined earlier and would have been kept; by the \omterm{def:b1:spanning-trees:cut}{cut} property, it belongs to a \omterm{def:b1:spanning-trees:mst}{minimum spanning tree}
together with the edges kept before it (the exchange of the proof never removes them, since none of them \omterm{def:b1:spanning-trees:cut}{crosses} that
\omterm{def:b1:spanning-trees:cut}{cut}). By induction, the \omterm{def:b1:graph-traversals:free-tree}{forest} is always contained in some \omterm{def:b1:spanning-trees:mst}{minimum spanning tree}, and at $n - 1$ edges it is one.

\omoutput{code/bachelor-1/26-spanning-trees/outputs/mst_town.txt}{The town as an undirected \omterm{def:b1:graph-traversals:graph}{graph}: fibre carries data both
ways, so the one-way road is an edge like the others.}\label{out:b1:spanning-trees:town}

\omcode{code/bachelor-1/26-spanning-trees/listings/kruskal.py}{8}{16}{Kruskal's algorithm, with \omterm{def:b1:spanning-trees:union-find}{union--find} for the
cycle test}\label{lst:b1:spanning-trees:kruskal}

\omoutput{code/bachelor-1/26-spanning-trees/outputs/kruskal.txt}{The fibre network of Brenmoor: 27.2 km of the 61 km
of road.}\label{out:b1:spanning-trees:kruskal}

\begin{omfigure}
\begin{tikzpicture}
\foreach \file/\title/\dx in {forest/{after 15 roads}/0, mst/{all 29 roads}/72} {
\begin{axis}[omrecord, xshift=\dx mm, width=7.2cm, height=6.6cm, xmin=-0.2, xmax=6, ymin=-0.2, ymax=5.1,
    axis equal image, grid=none, title={\small \title}, xticklabels={}, yticklabels={}]
  \addplot[black!20, line width=0.5pt] table[x=x, y=y]{code/bachelor-1/26-spanning-trees/figures/roads.dat};
  \addplot[omDef, line width=1.5pt] table[x=x, y=y]{code/bachelor-1/26-spanning-trees/figures/\file.dat};
  \addplot[only marks, mark=*, mark size=1.2pt, black!60] table[x=x, y=y]
    {code/bachelor-1/26-spanning-trees/figures/places.dat};
\end{axis}}
\end{tikzpicture}
\par\smallskip{\small Kruskal's algorithm on Brenmoor, drawn from the data by \code{figures/mst_map.py}: on the left a
\omterm{def:b1:graph-traversals:free-tree}{forest} of short roads, scattered over the town; on the right the \omterm{def:b1:spanning-trees:mst}{minimum spanning tree}.}
\end{omfigure}

Sorting costs $O(m \log m) = O(m \log n)$. The rest is $m$ tests ``same tree?'' and $n - 1$ merges; done naively, by
searching the \omterm{def:b1:graph-traversals:free-tree}{forest}, each test costs $O(n)$, and Kruskal's algorithm $O(mn)$.

\section{Disjoint sets}\label{sec:b1:spanning-trees:union-find}

\begin{definition}[Disjoint-set structure, representative]\label{def:b1:spanning-trees:union-find}
A \emph{disjoint-set structure}\index{disjoint-set structure}, or \emph{union--find}\index{union--find}, keeps a
partition of a set of elements, at first into singletons, under two operations: \code{find(x)} returns the
\emph{representative}\index{representative (of a set)} of the set of $x$, the same element for all members of a set;
\code{union(x, y)} merges the sets of $x$ and $y$.
\end{definition}

Two naive versions come first. One stores a label per element: \code{find} reads it in $O(1)$, \code{union} relabels a
whole set, $O(n)$. The other stores, for each element, a \emph{parent}; the \omterm{def:b1:spanning-trees:union-find}{representative} is the root, reached by
climbing: \code{union} links one root under the other, and \code{find} costs the height of the tree.

\omcodes{code/bachelor-1/26-spanning-trees/listings/uf_naive.py}{15-21,30-37}{The union of labels and the find of
parent \omterm{def:b1:pointers-and-arrays:pointer}{pointers}}\label{lst:b1:spanning-trees:naive}

\omoutput{code/bachelor-1/26-spanning-trees/outputs/uf_naive.txt}{A chain of $n - 1$ unions each followed by a find: both
versions are quadratic.}\label{out:b1:spanning-trees:naive}

\begin{definition}[Union by rank, path compression]\label{def:b1:spanning-trees:rank}
Each root carries a \emph{rank}\index{rank (union--find)}, 0 at the start. \emph{Union by rank}\index{union by rank}
links the root of smaller rank under the other, and raises the rank of the new root by one when the two ranks were
equal. \emph{Path compression}\index{path compression} makes every element met by a \code{find} point straight at the
root.
\end{definition}

\omcode{code/bachelor-1/26-spanning-trees/listings/union_find.py}{5}{30}{\omterm{def:b1:spanning-trees:rank}{Union by rank} and \omterm{def:b1:spanning-trees:rank}{path compression}, with the
parent steps counted}\label{lst:b1:spanning-trees:union-find}

\begin{proposition}[Height under union by rank]\label{prop:b1:spanning-trees:rank}
Without \omterm{def:b1:spanning-trees:rank}{path compression}, the tree of a root of \omterm{def:b1:spanning-trees:rank}{rank} $r$ has height $r$ and at least $2^r$ elements. So with $n$
elements every tree has height at most $\log_2 n$, and \code{find} costs $O(\log n)$.
\end{proposition}

\begin{proof}
A \omterm{def:b1:spanning-trees:rank}{rank} changes only when two trees of equal \omterm{def:b1:spanning-trees:rank}{rank} $r$ are joined: the new tree has height $r + 1$ and at least $2^r +
2^r$ elements. Joining under a root of larger \omterm{def:b1:spanning-trees:rank}{rank} changes neither its \omterm{def:b1:spanning-trees:rank}{rank} nor its height, which already exceeds the
other's by one at least. The size bound is \cref{exo:b1:spanning-trees:9}; then $2^r \leq n$ gives $r \leq \log_2 n$.
\end{proof}

\omoutput{code/bachelor-1/26-spanning-trees/outputs/rank_height.txt}{The tallest tree after random unions, and after
merging trees of equal \omterm{def:b1:spanning-trees:rank}{rank} two by two: the bound $\log_2 n$ is reached.}\label{out:b1:spanning-trees:height}

\begin{omfigure}
\begin{tikzpicture}[v/.style={circle, draw, minimum size=5mm, inner sep=0pt, font=\scriptsize}, >=stealth,
  level distance=9mm, sibling distance=9mm]
\node[v] (r) at (0,0) {r} child {node[v] (a) {a} child {node[v] (b) {b} child {node[v] (c) {c}}}} child {node[v]
{d}};
\node[font=\scriptsize] at (0.2,-3.7) {before \code{find(c)}};
\draw[->, omThm, thick] (1.3,-1.4) -- node[above, font=\scriptsize] {compress} (2.6,-1.4);
\begin{scope}[xshift=45mm, sibling distance=8mm]
\node[v] at (0,0) {r} child {node[v] {a}} child {node[v] {b}} child {node[v] {c}} child {node[v] {d}};
\node[font=\scriptsize] at (0,-3.7) {after: every node met points at the root};
\end{scope}
\end{tikzpicture}
\par\smallskip{\small \omterm{def:b1:spanning-trees:rank}{Path compression}. The find from \code{c} climbs three parent \omterm{def:b1:pointers-and-arrays:pointer}{pointers}; afterwards \code{c},
\code{b} and \code{a} hang from the root, and the next find from any of them climbs one. \omterm{def:b1:spanning-trees:rank}{Ranks} are not updated: a \omterm{def:b1:spanning-trees:rank}{rank}
remains an upper bound on the height.}
\end{omfigure}

With both improvements, a sequence of $m$ operations on $n$ elements costs $O(m\, \alpha(n))$, where $\alpha$ is an
inverse of Ackermann's function of \cref{prop:b1:induction-and-termination:ackermann}: it grows so slowly that
$\alpha(n) \leq 4$ for any $n$ that can be written down. Tarjan proved this bound in 1975; its proof is beyond this
book, and we admit it.

\omoutput{code/bachelor-1/26-spanning-trees/outputs/union_find.txt}{The chain of \cref{out:b1:spanning-trees:naive}
with both improvements: the work grows linearly.}\label{out:b1:spanning-trees:fast}

Kruskal's algorithm now costs $O(m \log n)$ for the sort and $O(m\, \alpha(n))$ for the rest: on the town, 65 parent
steps in all for the 53 roads.

\section{Prim's algorithm}\label{sec:b1:spanning-trees:prim}

\begin{method}[Prim's algorithm]\label{met:b1:spanning-trees:prim}
Grow a single tree from a start \omterm{def:b1:graph-traversals:graph}{vertex}: repeatedly add the lightest edge with exactly one end in the tree, and its
other end.
\end{method}

Each added edge is a lightest edge across the \omterm{def:b1:spanning-trees:cut}{cut} between the tree and the rest, so the \omterm{def:b1:spanning-trees:cut}{cut} property proves the
algorithm, as for Kruskal's. A heap holds the edges leaving the tree; an edge whose two ends have entered the tree is
discarded when it comes off the heap, as with lazy deletion in Dijkstra's algorithm.

\omcode{code/bachelor-1/26-spanning-trees/listings/prim.py}{9}{27}{Prim's algorithm with a heap of
edges}\label{lst:b1:spanning-trees:prim}

\omoutput{code/bachelor-1/26-spanning-trees/outputs/prim.txt}{From two different places, the same weight as
Kruskal's.}\label{out:b1:spanning-trees:prim}

Each edge enters the heap at most twice: $O(m \log n)$. Prim's algorithm suits dense \omterm{def:b1:graph-traversals:graph}{graphs}, given as a matrix: with an
\omterm{def:b1:pointers-and-arrays:array}{array} of best edges instead of a heap it costs $O(n^2)$, with no sort of the $m \approx n^2 / 2$ edges. Kruskal's
suits sparse \omterm{def:b1:graph-traversals:graph}{graphs} and edges that arrive as a list. Kruskal published his algorithm in 1956, and Prim his in 1957, for
networks of telephone lines.

\section{Bottlenecks}\label{sec:b1:spanning-trees:bottleneck}

\begin{definition}[Bottleneck spanning tree]\label{def:b1:spanning-trees:bottleneck}
A \emph{bottleneck spanning tree}\index{bottleneck spanning tree} is a \omterm{def:b1:spanning-trees:mst}{spanning tree} whose heaviest edge is as light as
possible among all \omterm{def:b1:spanning-trees:mst}{spanning trees}.
\end{definition}

A \omterm{def:b1:spanning-trees:mst}{minimum spanning tree} is a \omterm{def:b1:spanning-trees:bottleneck}{bottleneck spanning tree} (\cref{exo:b1:spanning-trees:11}). More: for any two \omterm{def:b1:graph-traversals:graph}{vertices},
the path joining them in a \omterm{def:b1:spanning-trees:mst}{minimum spanning tree} minimises the heaviest edge among all paths between them. Read with
maxima, a maximum \omterm{def:b1:spanning-trees:mst}{spanning tree} gives the \emph{widest} paths, those whose narrowest edge is widest, for every pair at
once. Removing the $k - 1$ heaviest edges of a \omterm{def:b1:spanning-trees:mst}{minimum spanning tree} splits the \omterm{def:b1:graph-traversals:graph}{vertices} into $k$ clusters, which
maximise the smallest distance between two clusters; the second volume uses this for clustering data.

\section{Exercises}\label{sec:b1:spanning-trees:exercises}

\begin{exercise}[$\star$]\label{exo:b1:spanning-trees:1}
A \omterm{def:b1:graph-traversals:graph}{graph} on $a, b, c, d, e$ has the edges $ab$ 4, $ac$ 1, $bc$ 3, $bd$ 2, $cd$ 5, $ce$ 7, $de$ 6. Give the edges kept by
Kruskal's algorithm, in order, and the weight.
\end{exercise}

\begin{exercise}[$\star$]\label{exo:b1:spanning-trees:2}
Run Prim's algorithm from $e$ on the same \omterm{def:b1:graph-traversals:graph}{graph}: in which order do the \omterm{def:b1:graph-traversals:graph}{vertices} enter?
\end{exercise}

\begin{exercise}[$\star$]\label{exo:b1:spanning-trees:3}
Starting from six singletons $0, \dots, 5$, run \code{union(0, 1)}, \code{union(2, 3)}, \code{union(1, 3)},
\code{union(4, 5)}, \code{union(5, 0)} with \cref{lst:b1:spanning-trees:union-find}. Give the parents and \omterm{def:b1:spanning-trees:rank}{ranks} at the
end.
\end{exercise}

\begin{exercise}[$\star$]\label{exo:b1:spanning-trees:4}
In the \omterm{def:b1:graph-traversals:graph}{graph} of \cref{exo:b1:spanning-trees:1}, is the tree $ac, bc, bd, de$ a \omterm{def:b1:spanning-trees:mst}{minimum spanning tree}? Use the cycle
property.
\end{exercise}

\begin{exercise}[$\star\star$]\label{exo:b1:spanning-trees:5}
This version replaces \omterm{def:b1:spanning-trees:union-find}{union--find} by a set of visited places. Why does it keep only 21 roads?
\omcode{code/bachelor-1/26-spanning-trees/listings/exo_visited.py}{6}{12}{Kruskal with a visited flag}
\omoutput[2-2]{code/bachelor-1/26-spanning-trees/outputs/exo_visited.txt}{On the town.}
\end{exercise}

\begin{exercise}[$\star\star$]\label{exo:b1:spanning-trees:6}
Find the error in this \omterm{def:b1:spanning-trees:rank}{union by rank}, and explain the tree of height 999.
\omcode{code/bachelor-1/26-spanning-trees/listings/exo_rank_bug.py}{7}{16}{A wrong \omterm{def:b1:spanning-trees:rank}{rank}}
\omoutput[2-3]{code/bachelor-1/26-spanning-trees/outputs/exo_rank_bug.txt}{The unions \code{union(i, i - 1)} for $i$ from 1
to 999, before any other find.}
\end{exercise}

\begin{exercise}[$\star\star$]\label{exo:b1:spanning-trees:7}
Compute the weight of a second-best \omterm{def:b1:spanning-trees:mst}{spanning tree}, the lightest \omterm{def:b1:spanning-trees:mst}{spanning tree} different from the minimum one, by
swapping one edge of the \omterm{def:b1:spanning-trees:mst}{minimum spanning tree} for one outside it; test it against all \omterm{def:b1:spanning-trees:mst}{spanning trees} of small
\omterm{def:b1:graph-traversals:graph}{graphs}.
\end{exercise}

\begin{exercise}[$\star\star$]\label{exo:b1:spanning-trees:8}
Generate a maze on a grid of cells by Kruskal's algorithm on the walls taken in a random order, seeded. Why is every
cell reachable from every other, by exactly one way?
\end{exercise}

\begin{exercise}[$\star\star$]\label{exo:b1:spanning-trees:9}
Show that with \omterm{def:b1:spanning-trees:rank}{union by rank} a root of \omterm{def:b1:spanning-trees:rank}{rank} $r$ has at least $2^r$ elements in its tree, with or without \omterm{def:b1:spanning-trees:rank}{path
compression}.
\end{exercise}

\begin{exercise}[$\star\star\star$]\label{exo:b1:spanning-trees:10}
Show that a \omterm{def:b1:graph-traversals:connected}{connected graph} whose edge weights are all distinct has exactly one \omterm{def:b1:spanning-trees:mst}{minimum spanning tree}.
\end{exercise}

\begin{exercise}[$\star\star\star$]\label{exo:b1:spanning-trees:11}
Show that every \omterm{def:b1:spanning-trees:mst}{minimum spanning tree} is a \omterm{def:b1:spanning-trees:bottleneck}{bottleneck spanning tree}, and give a \omterm{def:b1:spanning-trees:bottleneck}{bottleneck spanning tree} that is not
minimum.
\end{exercise}

\begin{exercise}[$\star\star\star$]\label{exo:b1:spanning-trees:12}
Borůvka's method works in rounds: each tree of the \omterm{def:b1:graph-traversals:free-tree}{forest}, at first the single \omterm{def:b1:graph-traversals:graph}{vertices}, picks its lightest leaving
edge, and all picked edges are added. Implement it with \omterm{def:b1:spanning-trees:union-find}{union--find}, prove that it needs at most $\log_2 n$ rounds,
and test it against Kruskal's algorithm.
\end{exercise}

\section{Problem: Laying fibre in Brenmoor}\label{sec:b1:spanning-trees:problem}

\begin{problem}[{Weekend problem --- Laying fibre in Brenmoor}]\label{pb:b1:spanning-trees:1}
The council has chosen to lay fibre along a \omterm{def:b1:spanning-trees:mst}{minimum spanning tree} of the roads, weighted by length. Before trusting the
algorithms with the whole town, it counts the trees of a small piece of it.
\omcode{code/bachelor-1/26-spanning-trees/listings/count_trees.py}{11}{16}{Every \omterm{def:b1:spanning-trees:mst}{spanning tree}, by \omterm{def:b1:backtracking:exhaustive}{brute
force}}\label{lst:b1:spanning-trees:count}
\omoutput{code/bachelor-1/26-spanning-trees/outputs/count_trees.txt}{The places within two roads of the Town
Square.}\label{out:b1:spanning-trees:count}
\omoutput{code/bachelor-1/26-spanning-trees/outputs/bottleneck.txt}{Minimax routes read from the tree, and clusters
obtained by removing its longest roads.}\label{out:b1:spanning-trees:bottleneck}

\medskip\noindent\textbf{Part I --- Spanning trees of a piece of the town.}
\begin{enumerate}
\item Why does a set of $n - 1$ edges without a cycle form a \omterm{def:b1:spanning-trees:mst}{spanning tree}?
\item Why are only 37 of the 165 sets of 8 roads \omterm{def:b1:spanning-trees:mst}{spanning trees}?
\item How many sets would the \omterm{def:b1:backtracking:exhaustive}{brute force} try for the whole town, $\binom{53}{29}$? Estimate its order of magnitude.
\item Why does \code{all} stop at the first union that fails, and why is that correct?
\item The lightest tree found by \omterm{def:b1:backtracking:exhaustive}{brute force} weighs the same as Kruskal's: what does that check, and what not?
\item How many cycles at least does the piece have, given 9 places and 11 roads?
\end{enumerate}

\medskip\noindent\textbf{Part II --- The whole town.}
\begin{enumerate}[resume]
\item Read the fibre length and the number of roads used.
\item Why do Kruskal's and Prim's algorithms give the same weight from any start?
\item Are the road lengths distinct? If not, could two \omterm{def:b1:spanning-trees:mst}{minimum spanning trees} exist?
\item Count the heap pops of Prim's algorithm, and explain the 44.
\item Which of the two algorithms would you use for a town of a million places and three million roads?
\item Why is the longest road of the tree, 1540 m, kept although longer roads exist?
\end{enumerate}

\medskip\noindent\textbf{Part III --- Union--find.}
\begin{enumerate}[resume]
\item Explain the 9900 and 9801 steps of \cref{out:b1:spanning-trees:naive}.
\item Why does the chain cost about $2n$ steps with \omterm{def:b1:spanning-trees:rank}{union by rank} and \omterm{def:b1:spanning-trees:rank}{path compression}?
\item Why does the largest \omterm{def:b1:spanning-trees:rank}{rank} stay at 1 on this chain?
\item Why do random unions give trees much lower than $\log_2 n$?
\item Describe the unions that reach height $\log_2 n$ exactly.
\item Why does \code{find} not update \omterm{def:b1:spanning-trees:rank}{ranks} when it compresses?
\end{enumerate}

\medskip\noindent\textbf{Part IV --- Bottlenecks.}
\begin{enumerate}[resume]
\item Prove that for two \omterm{def:b1:graph-traversals:graph}{vertices} $a$ and $b$, the path from $a$ to $b$ in a \omterm{def:b1:spanning-trees:mst}{minimum spanning tree} has the smallest
possible heaviest edge among all paths from $a$ to $b$.
\item What does ``no route avoids a road of 1240 m'' mean for a lorry that must avoid long roads?
\item How was the bottleneck computed without the tree, and why do the two answers agree?
\item Explain the clusters of sizes 29 and 1, then 26, 3 and 1.
\item Is removing the longest tree roads a good way to split the town into districts?
\item State the named result, \emph{the minimum fibre length of Brenmoor}, with the bottleneck theorem.
\end{enumerate}
\end{problem}
